\subsection{P005 --- CO2 concentration prediction in office spaces using physics-informed neural network based on number of occupants and IoT sensor data}
\label{paper:P005}
\textbf{Authors:} Jehyun Kim, Gihoon Kim, Jong-Il Bang, Anseop Choi, Minki Sung\par
\textbf{Major Category:} Building Environment and IEQ\par
\textbf{Secondary Category:} AI and Machine Learning for Buildings, Occupancy and Human-Building Interaction, Sensors, IoT and Data Integration\par
\textbf{Keywords:} CO2 prediction, Physics-informed neural network, Occupancy, HVAC operation, IoT, Indoor air quality\par
\paragraph{主な主張}
occupant count, AHU operating status, outdoor-air damper opening, outdoor air quality等を入力し, CO2 mass balance equationをlossに組み込んだPINNを評価した. 9th floor dataで10-min予測RMSE 4.04 ppm, 60-min recursive予測RMSE 17.16 ppmを報告した. \cite{Kim:2026CO2PINN}
\paragraph{新規性}
CO2 mass balance equationをlearning objectiveへ直接組み込み, 実測occupancyとHVAC operational dataを同時利用することで, data-driven predictionへ物理制約を付加した点に特徴がある.
\paragraph{位置付け}
office CO2 predictionとIAQが主要評価対象であるためBuilding Environment and IEQを主分類とし, PINN, occupancy, IoT/HVAC data利用を副分類で表現する.
